\chapter{Configuración experimental y resultados}\label{cap:resultados}

En este capítulo se presentan la configuración experimental y los
resultados obtenidos al entrenar y evaluar los agentes descritos en el
Capítulo~\ref{cap:diseno-arquitectura}. Primero se detallan los
escenarios de la rejilla, los hiperparámetros de los algoritmos, el
entorno de ejecución y el protocolo de evaluación. A continuación se
presentan los resultados del agente PPO de referencia (test 25),
seguidos de una serie de estudios complementarios que modifican un único
elemento del baseline para medir su efecto: aprendizaje por currículo,
reward geodésico, política recurrente (GRU), alcance del LiDAR y
enmascarado de acciones inválidas. Al final del capítulo se discuten los
patrones observados.

\section{Configuración experimental}\label{sec:config_experimental}

\subsection{Escenarios de la rejilla}
Los experimentos se realizan sobre cuatro mapas estáticos de ocupación,
almacenados como ficheros JSON y cargados en tiempo de ejecución.
Durante el entrenamiento, el mapa se selecciona de manera uniforme al
inicio de cada episodio (variante multi-mapa del entorno
\texttt{Grid\_Robot\_Sections\_Env\_MultiMap}); durante la evaluación, el
protocolo se describe en la Sección~\ref{sec:protocolo}. Los mapas
empleados se resumen en el Cuadro~\ref{tab:mapas}.

\begin{table}[htbp]
    \centering
    \caption{Mapas de ocupación empleados en los experimentos. Las
    celdas ocupadas incluyen tanto los obstáculos declarados como su
    posible solapamiento.}
    \label{tab:mapas}
    \begin{tabular}{llc}
        \toprule
        \textbf{Mapa} & \textbf{Descripción} & \textbf{Celdas ocupadas} \\
        \midrule
        \texttt{map\_obstacle\_free}  & Rejilla vacía, caso de control. & $0$ \\
        \texttt{map\_1\_obstacle}      & Un obstáculo asimétrico. & $\sim 20$ \\
        \texttt{map\_4\_square\_obstacle} & Cuatro obstáculos rectangulares. & $\sim 60$ \\
        \texttt{map\_simple}           & 14 obstáculos solapados que forman un bolsillo cóncavo. & $121$ \\
        \bottomrule
    \end{tabular}
\end{table}

El escenario \texttt{map\_obstacle\_free} valida que la política aprenda
el término de progreso aislado del resto de términos de la recompensa.
Los escenarios con obstáculos aumentan progresivamente la densidad
geométrica; \texttt{map\_simple} es el más exigente porque contiene una
geometría cóncava en la que la política reactiva debe alejarse
momentáneamente de la meta para rodear el bolsillo, lo que configura
precisamente el modo de fallo por mínimo local descrito por
\citet{gao2025hierarchical}.

\subsection{Hiperparámetros del agente PPO}
El experimento de referencia (test 25) entrena el agente
\texttt{PPOAgent} con los hiperparámetros del
Cuadro~\ref{tab:hiperparametros_entrenamiento}. Todos los experimentos
posteriores (tests 26--31) heredan esta misma configuración y modifican
un único elemento, descrito en cada apartado.

\begin{table}[htbp]
    \centering
    \caption{Hiperparámetros del agente PPO en el experimento de
    referencia (test 25).}
    \label{tab:hiperparametros_entrenamiento}
    \begin{tabular}{lll}
        \toprule
        \textbf{Hiperparámetro} & \textbf{Valor} & \textbf{Descripción} \\
        \midrule
        Tasa de aprendizaje $\alpha_{lr}$ & $10^{-3}$ & Optimizador Adam, igual para actor y crítico. \\
        Factor de descuento $\gamma$ & $0{,}99$ & Descuento temporal del retorno. \\
        $\lambda$ de GAE & $0{,}95$ & Suavizado del estimador de ventajas. \\
        Clip PPO $\varepsilon$ & $0{,}20$ & Recorte del objetivo \textit{surrogate}. \\
        Coef.\ de valor $c_v$ & $0{,}50$ & Peso de la pérdida del crítico. \\
        Coef.\ de entropía $c_H$ & $0{,}01$ & Peso del término de entropía. \\
        Tamaño de lote & $64$ & Mini-lote del optimizador. \\
        Capacidad del \textit{buffer} & $10^{5}$ & \textit{Buffer} circular de transiciones. \\
        Norma máx.\ de gradiente & $5{,}0$ & \textit{Gradient clipping} aplicado al optimizador. \\
        Episodios de entrenamiento & $40\,000$ & Horizonte total. \\
        Horizonte por episodio & $400$ pasos & Trunca el episodio si no se alcanza la meta antes. \\
        Pasos por actualización & $4$ & \texttt{TrainerConfig.eps\_step}, actualiza cada 4 pasos. \\
        \bottomrule
    \end{tabular}
\end{table}

\subsection{Hiperparámetros del agente DQN}
Para el agente DQN basado en \texttt{GridAgent} se emplean los
hiperparámetros del Cuadro~\ref{tab:hiperparametros_dqn}, utilizando la
variante Double DQN con red objetivo congelada.

\begin{table}[htbp]
    \centering
    \caption{Hiperparámetros del agente DQN (\texttt{GridAgent}).}
    \label{tab:hiperparametros_dqn}
    \begin{tabular}{lll}
        \toprule
        \textbf{Hiperparámetro} & \textbf{Valor} & \textbf{Descripción} \\
        \midrule
        Tasa de aprendizaje $\alpha_{lr}$ & $10^{-3}$ & Optimizador Adam. \\
        Factor de descuento $\gamma$ & $0{,}99$ & Descuento temporal del retorno. \\
        $\varepsilon_{\max}$ & $1{,}00$ & Probabilidad inicial de exploración. \\
        $\varepsilon_{\min}$ & $0{,}05$ & Probabilidad asintótica de exploración. \\
        $\tau_{\varepsilon}$ & $10^{5}$ pasos & Constante de decaimiento exponencial de $\varepsilon$. \\
        Frec.\ actualización red objetivo & $10^{3}$ pasos & Sincroniza $\theta^- \leftarrow \theta$. \\
        Tamaño de lote & $64$ & Mini-lote del optimizador. \\
        Capacidad del \textit{buffer} & $10^{5}$ & \textit{Buffer} circular de transiciones. \\
        Norma máx.\ de gradiente & $5{,}0$ & \textit{Gradient clipping}. \\
        Variante & Double DQN & Desacopla selección y evaluación de acciones. \\
        \bottomrule
    \end{tabular}
\end{table}

\subsection{Función de recompensa}
Todos los experimentos emplean la función
\texttt{ProgressReward} descrita en la
Sección~\ref{sec:recompensa}, con los parámetros del
Cuadro~\ref{tab:recompensa_params} salvo donde se indique lo contrario
(test 27 activa la variante geodésica).

\begin{table}[htbp]
    \centering
    \caption{Parámetros de la función de recompensa
    \texttt{ProgressReward} en el experimento de referencia.}
    \label{tab:recompensa_params}
    \begin{tabular}{ll}
        \toprule
        \textbf{Parámetro} & \textbf{Valor} \\
        \midrule
        Escala de progreso $w_d$ & $10$ \\
        Coste por paso $c_{\text{step}}$ & $0{,}05$ \\
        Margen de seguridad $D_{\text{safe}}$ & $\tfrac{2}{3}\,d_{\max}$ ($\approx 6{,}67$ celdas) \\
        Penalización máx.\ por proximidad $\lambda_{\text{prox}}$ & $0{,}01$ \\
        Recompensa terminal por meta $R_{\text{win}}$ & $+20$ \\
        Penalización por colisión $R_{\text{fail}}$ & $-5$ \\
        Tolerancia a la meta $\varepsilon_{\text{goal}}$ & $1$ celda \\
        Progreso negativo recortado a cero & Sí \\
        Métrica de distancia & Euclídea (geodésica en test 27) \\
        \bottomrule
    \end{tabular}
\end{table}

\subsection{Configuración del sensor LiDAR}
El sensor \texttt{GridLidar} se configura con alcance máximo
$d_{\max} = 10$ celdas, campo de visión $2\pi$ y ruido gaussiano
aditivo $\sigma_{\text{noise}} = 10^{-3}$ celdas. El número total de
rayos es $N_{\text{rays}} = 15\,N_{\text{sec}}$ para garantizar varias
muestras por sector; en el experimento de referencia
$N_{\text{sec}} = 8$ (y por tanto $N_{\text{rays}} = 120$). El sensor
mantiene una caché de distancias por mapa, acelerada en GPU, que se
reutiliza entre episodios siempre que el campo de visión sea $2\pi$.

\subsection{Entorno de ejecución}
El entrenamiento se ejecuta en el equipo del autor con una tarjeta
NVIDIA compatible con CUDA; cuando no hay GPU disponible, el agente
y el sensor delegan en la CPU (en MPS si se detecta un Apple Silicon).
El tiempo total del test 25 fue de 5\,h 49\,min para $40\,000$
episodios en GPU, lo que da una velocidad media aproximada de $115$
episodios por minuto. La capa de simulación emplea
\texttt{gymnasium} 0.29+ y PyTorch; el \textit{logging} de métricas se
realiza con \texttt{TensorBoard} y la configuración de cada ejecución se
serializa en un fichero \texttt{config.json} dentro del directorio
\texttt{runs/Test/<N>/}.

\subsection{Protocolo de evaluación}\label{sec:protocolo}
Al final del entrenamiento, el agente se evalúa durante $1\,000$
episodios sobre los cuatro mapas, con la semilla $0$ y en dos
modalidades complementarias:
\begin{itemize}
    \item \textbf{Modo \textit{sampled}}: la acción se muestrea de la
    distribución categórica del actor, replicando el comportamiento
    estocástico del entrenamiento.
    \item \textbf{Modo \textit{greedy}}: la acción se escoge
    deterministamente como el $\arg\max$ de los \textit{logits} del
    actor; refleja el comportamiento más conservador que cabe esperar
    en despliegue.
\end{itemize}

Se reportan las siguientes métricas globales y por mapa:
\begin{itemize}
    \item Tasa de éxito (\textit{Success rate}, SR): porcentaje de
    episodios terminados por alcanzar la meta.
    \item Tasa de colisión y tasa de agotamiento del horizonte temporal.
    \item \textit{Loops}: número de episodios en los que la trayectoria
    visita alguna celda al menos veinte veces, indicador de ciclos
    estables (típicos de la modalidad \textit{greedy}).
    \item \textit{Path ratio}: cociente entre los pasos efectivos del
    episodio y el camino óptimo calculado por BFS; valores próximos a
    uno indican trayectorias eficientes.
\end{itemize}

Las comparaciones entre experimentos emplean el test $z$ de proporciones
de dos muestras ($|z|>1{,}96$ equivale a significación estadística al
$5\,\%$).

\section{Resultados del baseline PPO (test 25)}\label{sec:resultados_baseline}
El agente PPO del baseline alcanza una tasa de éxito acumulada del
$83\,\%$ en los primeros $4\,000$ episodios y asciende hasta el
$93\,\%$ al llegar a los $40\,000$ episodios, con la entropía de la
política todavía en $0{,}55$ nats al final del entrenamiento. En la
evaluación posterior (1\,000 episodios, semilla $0$), los resultados
globales se recogen en el Cuadro~\ref{tab:baseline_global}.

\begin{table}[htbp]
    \centering
    \caption{Resultados del agente PPO de referencia (test 25) en la
    evaluación final sobre los cuatro mapas (1\,000 episodios por
    modalidad, semilla $0$). Los \textit{loops} se cuentan sobre los
    episodios que agotan el horizonte temporal.}
    \label{tab:baseline_global}
    \begin{tabular}{lcccc}
        \toprule
        \textbf{Modalidad} & \textbf{SR global (\%)} & \textbf{Colisiones (\%)} & \textbf{Timeouts (\%)} & \textbf{Loops} \\
        \midrule
        Sampled & $93{,}9$ & $5{,}1$ & $1{,}0$ & $3$ \\
        Greedy  & $91{,}1$ & $1{,}1$ & $7{,}8$ & $76$ \\
        \bottomrule
    \end{tabular}
\end{table}

En la evaluación por mapa (Cuadro~\ref{tab:baseline_por_mapa}), tres de
los cuatro escenarios superan el $91\,\%$ de éxito en ambas modalidades,
mientras que \texttt{map\_simple} se estabiliza en torno al $75$--$78\,\%$
según la modalidad. La diferencia entre modalidades es especialmente
visible en \texttt{map\_4\_square\_obstacle}, que pasa de $91{,}4\,\%$ en
\textit{greedy} a $99{,}2\,\%$ al muestrear la política. Dicho de otro
modo: el modo \textit{sampled} cambia los \textit{timeouts} del
\textit{greedy} por un número mayor de colisiones, lo que se traduce en
una tasa de éxito global más alta ($93{,}9\,\%$ frente a $91{,}1\,\%$)
pero a costa de más pasos cerca de los obstáculos.

El análisis de las trayectorias fallidas muestra que la mayor parte de
las colisiones en \texttt{map\_simple} se producen con el LiDAR
detectando el obstáculo a una sola celda de distancia en la dirección de
avance: no es un problema de percepción, sino de ruido residual de la
política junto a una pared, que en la modalidad \textit{greedy} degenera
en ciclos estables ($76$ de los $78$ \textit{timeouts} son
\textit{loops}). Las trayectorias que sí alcanzan la meta son
prácticamente óptimas: la mediana del cociente \textit{pasos\,/\,BFS}
vale $1{,}0$ en \textit{greedy} para los cuatro mapas y oscila entre
$1{,}0$ y $1{,}1$ en \textit{sampled}.

\begin{table}[htbp]
    \centering
    \caption{Tasa de éxito por mapa del baseline PPO (test 25). Entre
    paréntesis, el número de episodios con el mapa seleccionado sobre el
    total de $1\,000$.}
    \label{tab:baseline_por_mapa}
    \begin{tabular}{lccc}
        \toprule
        \textbf{Mapa} & \textbf{Episodios} & \textbf{SR \textit{sampled} (\%)} & \textbf{SR \textit{greedy} (\%)} \\
        \midrule
        \texttt{map\_obstacle\_free}  & $267$ & $99{,}6$ & $99{,}6$ \\
        \texttt{map\_1\_obstacle}      & $230$ & $98{,}7$ & $98{,}7$ \\
        \texttt{map\_4\_square\_obstacle} & $255$ & $99{,}2$ & $91{,}4$ \\
        \texttt{map\_simple}           & $248$ & $77{,}8$ & $74{,}6$ \\
        \midrule
        \textbf{Global}                & $1\,000$ & $\mathbf{93{,}9}$ & $\mathbf{91{,}1}$ \\
        \bottomrule
    \end{tabular}
\end{table}

\section{Estudios complementarios}\label{sec:estudios_complementarios}
Los experimentos siguientes modifican un único elemento del baseline
para aislar su efecto. El Cuadro~\ref{tab:comparativa_experimentos}
resume los resultados globales; los apartados posteriores detallan cada
estudio.

\begin{table}[htbp]
    \centering
    \caption{Comparación de los experimentos complementarios frente al
    baseline PPO (test 25). Las columnas $z$ reportan el estadístico
    del test de proporciones frente al baseline sobre $1\,000$ episodios
    por modalidad; $|z| > 1{,}96$ equivale a significación al $5\,\%$.}
    \label{tab:comparativa_experimentos}
    \begin{tabular}{lccccc}
        \toprule
        \textbf{Test} & \textbf{Variante} & \textbf{SR \textit{sampled} (\%)} & \textbf{SR \textit{greedy} (\%)} & \textbf{SR map\_simple (\%)} & \textbf{Tiempo} \\
        \midrule
        25 & Baseline PPO                  & $93{,}9$ & $91{,}1$ & $77{,}8$ & 5\,h 49\,min \\
        26 & + currículo (5 etapas)        & $93{,}9$ ($z\!=\!0$) & $89{,}3$ (n.s.) & $77{,}8$ (n.s.) & 3\,h 18\,min \\
        27 & reward geodésico              & $93{,}4$ (n.s.) & $85{,}0$ ($z\!=\!-4{,}2$) & $78{,}6$ (n.s.) & 2\,h 52\,min \\
        28 & GRU, sin \textit{grad clip}   & $85{,}2$ & $77{,}4$ & $\sim 60$ & 4\,h 37\,min \\
        29 & GRU, con \textit{grad clip}   & $89{,}8$ ($z\!=\!-3{,}4$) & $88{,}3$ ($z\!=\!-2{,}1$) & $\sim 70$ & 2\,h 51\,min \\
        30 & LiDAR $d_{\max} = 30$         & $93{,}7$ (n.s.) & $90{,}4$ (n.s.) & $79{,}4$ (n.s.) & 2\,h 41\,min \\
        31 & + \textit{action masking}     & \multicolumn{3}{c}{\textit{pendiente de entrenamiento}} & -- \\
        \bottomrule
    \end{tabular}
\end{table}

\subsection{Aprendizaje por currículo (test 26)}
El test 26 añade al baseline un currículo \texttt{SuccessCurriculum} de
cinco etapas que empieza en el mapa vacío y añade progresivamente los
mapas con obstáculos, con un umbral de éxito del $90\,\%$ (etapas 1--4)
y del $85\,\%$ (etapa final) medido sobre una ventana de $500$
episodios. Las etapas se activan en los episodios
$0 / 1620 / 2120 / 3416 / 3916$, de modo que el agente llega al
escenario completo tras unos $4\,000$ episodios. El currículo reduce
sensiblemente el tiempo necesario para alcanzar el $80\,\%$ de éxito,
pero la política final es estadísticamente indistinguible del baseline
(SR \textit{sampled} $93{,}9\,\%$ en ambos). El resultado concuerda con
los informes de \citet{xue2022using}, en los que la ganancia del
currículo se concentra en la fase inicial del entrenamiento.

\subsection{Progreso geodésico (test 27)}
El test 27 sustituye la distancia euclídea del término de progreso por
la distancia geodésica (BFS 4-conexo desde la celda meta, respetando
los obstáculos). La hipótesis era que señalar al agente el camino que
rodea el obstáculo mejoraría el rendimiento en \texttt{map\_simple}.
En la práctica, el modo \textit{sampled} no mejora de forma
significativa (SR $93{,}4\,\%$ frente a $93{,}9\,\%$) y el modo
\textit{greedy} empeora de manera clara (SR $85{,}0\,\%$ frente a
$91{,}1\,\%$, $z = -4{,}2$) con un aumento notable de los \textit{loops}
($137$ frente a $76$). La señal geodésica parece inducir políticas
determinísticas que explotan trayectorias de bordes y quedan atrapadas
en ciclos cuando se evalúan en modo \textit{greedy}.

\subsection{Política recurrente con GRU (tests 28 y 29)}
Los tests 28 y 29 evalúan la variante \texttt{RecurrentPPOAgent}
descrita en la Sección~\ref{sec:ppo_lstm}, con una capa GRU de $40$
unidades y secuencias de longitud $32$. El test 28 reutilizó el bucle
de entrenamiento del PPO estándar y descubrió que la clase recurrente
no aplicaba el recorte de gradiente ni la máscara de padding; esto
produjo un colapso del rendimiento desde el $90\,\%$ hasta el $25\,\%$
alrededor del episodio $25\,000$. El test 29 corrige ambos problemas y
reduce la tasa de aprendizaje a $3\!\cdot\!10^{-4}$. El entrenamiento es
entonces estable, pero la política final queda por debajo del baseline
MLP (SR \textit{sampled} $89{,}8\,\%$, $z = -3{,}4$) con una entropía
mucho menor ($0{,}27$ nats) y entre $95$ y $111$ \textit{loops} en ambas
modalidades. Las trayectorias exitosas son casi óptimas (mediana del
\textit{path ratio} $= 1{,}00$), pero la memoria recurrente, tal como
está implementada, no elimina los ciclos. Este estudio no permite
concluir que la memoria no aporte valor en general; sí indica que, con
esta implementación y presupuesto de entrenamiento, el MLP sigue
siendo competitivo.

\subsection{Alcance extendido del LiDAR (test 30)}
El test 30 aumenta el alcance del sensor a $d_{\max} = 30$ celdas (frente
a $10$ en el baseline), manteniendo el resto de parámetros constantes.
El margen de seguridad relativo al nuevo alcance se recalcula para que
siga valiendo $6{,}67$ celdas en términos absolutos. El resultado es
estadísticamente indistinguible del baseline en todas las métricas
(SR \textit{sampled} $93{,}7\,\%$, $z = -0{,}2$), lo que sugiere que la
limitación principal del baseline no reside en la percepción.

\subsection{Enmascarado de acciones inválidas (test 31)}
El test 31 activa el enmascarado descrito en la
Sección~\ref{sec:action_masking} manteniendo el resto de la
configuración del baseline. En una prueba de humo de $40$ episodios no
se registraron colisiones. El experimento está preparado para
entrenamiento completo; los resultados definitivos se incorporarán a
este capítulo cuando esté disponible la ejecución de $40\,000$
episodios. Se espera que el enmascarado reduzca a cero las colisiones y
no modifique de manera sustancial el comportamiento en los episodios de
mínimo local, en los que ninguna acción individual lleva al choque y, por
tanto, la máscara no restringe el espacio efectivo de decisión.

\section{Discusión}\label{sec:discusion}
Los resultados experimentales permiten extraer varias observaciones
sobre los límites del enfoque adoptado y sobre la relación entre las
decisiones de diseño del Capítulo~\ref{cap:diseno-arquitectura} y el
comportamiento observado.

\textbf{Casi todas las variantes convergen a un rendimiento similar.}
Los tests 25 (baseline), 26 (currículo), 27 (reward geodésico) y 30
(LiDAR extendido) terminan en una ventana estrecha de tasa de éxito
global ($\sim 93\,\%$ en modo \textit{sampled}). Las variantes
recurrentes (28 y 29) son algo peores. La convergencia a un
rendimiento similar sugiere que, con el presupuesto de interacción
estudiado, el cuello de botella no se encuentra en el aproximador ni en
la percepción, sino en la señal de aprendizaje y en los patrones
geométricos a los que la política no sabe responder.

\textbf{El modo de fallo persistente es un mínimo local geométrico.}
\texttt{map\_simple} se estabiliza en torno al $78\,\%$ en todas las
variantes probadas. Las trayectorias fallidas siguen el patrón descrito
por \citet{gao2025hierarchical}: el robot queda detenido o entra en un
ciclo frente a un bolsillo cóncavo porque la salida requiere alejarse
transitoriamente de la meta, acción que la señal densa de progreso
desincentiva. Ni el currículo (que ajusta la distribución de tareas
pero no modifica la función de recompensa) ni la señal geodésica (que
dirige la política hacia la salida pero genera ciclos deterministas en
\textit{greedy}) resuelven este modo de fallo, y la memoria GRU
implementada en los tests 28--29 tampoco lo hace.

\textbf{La modalidad de evaluación cambia la lectura.} El modo
\textit{greedy} elimina la variabilidad útil que, en modo
\textit{sampled}, saca al agente de los ciclos: el baseline mejora de
$91{,}1\,\%$ a $93{,}9\,\%$ al pasar de \textit{greedy} a
\textit{sampled}, y el test 27 muestra la mayor brecha
($85{,}0\,\%$ frente a $93{,}4\,\%$). Esta diferencia es relevante para
una futura evaluación en un sistema real, donde una política demasiado
determinista puede quedar atrapada en \textit{loops} que no aparecerían
con un muestreo estocástico.

\textbf{Compromiso entre reutilización de muestras y estabilidad.} La
implementación de PPO empleada (Sección~\ref{sec:redes}) muestrea
mini-lotes de un \textit{buffer} con transiciones de episodios
anteriores, acercándose a un algoritmo \textit{off-policy}. Esta
decisión permite al agente aprovechar transiciones recolectadas con
versiones ligeramente antiguas de la política, con el coste de que los
log-probabilidades utilizados en la actualización pueden quedar
obsoletos. En el baseline el efecto no bloquea el aprendizaje, pero en
las variantes recurrentes se combina con la ausencia inicial de recorte
de gradiente para producir el colapso observado en el test 28. Una
línea de trabajo futuro inmediata consiste en comparar esta
implementación con un PPO estrictamente \textit{on-policy} sobre el
mismo escenario, para medir la contribución específica del
\textit{buffer}.

\textbf{Validez interna y limitaciones.} Las conclusiones de este
capítulo se circunscriben al robot \texttt{GridRobotNotTurning}, a los
cuatro mapas descritos, al algoritmo PPO con los hiperparámetros
documentados y a los $40\,000$ episodios de entrenamiento. Las
cifras reportadas corresponden a una única semilla ($0$); una
estimación de la varianza entre semillas, así como la comparación
sistemática con el agente DQN de la Sección~\ref{sec:redes}, quedan
como líneas de trabajo pendientes antes de la defensa.
